\section*{Capítulo \ref{ch:g11:circuits-and-power} --- Circuitos eléctricos y potencia}

\begin{solution}{exo:g11:circuits-and-power:1}
$q = It = 2.0 \times 90 \times 60 = \qty{1.1e4}{C}$. Número de cargas
elementales:
$q/e = 1.08\times10^{4}/1.6\times10^{-19} \approx \num{6.8e22}$.
\end{solution}

\begin{solution}{exo:g11:circuits-and-power:2}
$I = U/R = 12/220 = \qty{0.055}{A} = \qty{55}{mA}$;
$R = U/I = 230/0.50 = \qty{460}{\ohm}$.
\end{solution}

\begin{solution}{exo:g11:circuits-and-power:3}
En serie: $120 + 60 = \qty{180}{\ohm}$. En paralelo:
$\frac{120 \times 60}{180} = \qty{40}{\ohm}$ --- menor que las dos:
la misma \omterm{def:g11:circuits-and-power:voltage}{tensión}, pero la corriente dispone de un segundo camino, así
que hay más intensidad total y menos \omterm{def:g11:circuits-and-power:resistance}{resistencia}.
\end{solution}

\begin{solution}{exo:g11:circuits-and-power:4}
$P = UI = 230 \times 8.7 = \qty{2.0}{kW}$. En $t = \qty{240}{s}$:
$E = Pt = 2000 \times 240 = \qty{4.8e5}{J} = 4.8\times10^{5}/3.6
\times10^{6} \approx \qty{0.13}{kW.h}$.
\end{solution}

\begin{solution}{exo:g11:circuits-and-power:5}
$E = 60 \times 5.0 = \qty{300}{W.h} = \qty{0.30}{kW.h} = 0.30 \times
3.6\times10^{6} = \qty{1.1e6}{J}$. Coste: $0.30 \times 0.25 = 0.075$
euros al día.
\end{solution}

\begin{solution}{exo:g11:circuits-and-power:6}
\emph{1.} $R = U^2/P = 230^2/1200 = \qty{44}{\ohm}$;
$I = P/U = 1200/230 = \qty{5.2}{A}$.

\emph{2.} $P' = U'^2/R = 115^2/44.1 = \qty{300}{W}$: la cuarta parte, no
la mitad, porque reducir $U$ a la mitad reduce también $I$ a la mitad, y
$P = UI \propto U^2$.
\end{solution}

\begin{solution}{exo:g11:circuits-and-power:7}
$R_{\text{s}} = \qty{60}{\ohm}$, luego $I = 12/60 = \qty{0.20}{A}$.
$U_1 = R_1 I = \qty{8.0}{V}$, $U_2 = \qty{4.0}{V}$ (suman \qty{12}{V});
$P_1 = U_1 I = \qty{1.6}{W}$, $P_2 = \qty{0.8}{W}$: en total
\qty{2.4}{W} $= UI = 12 \times 0.20$, como debía ser.
\end{solution}

\begin{solution}{exo:g11:circuits-and-power:8}
$I_1 = 12/60 = \qty{0.20}{A}$, $I_2 = 12/30 = \qty{0.40}{A}$, en total
$I = \qty{0.60}{A}$. Equivalente: $R = U/I = 12/0.60 = \qty{20}{\ohm}$,
y en efecto $\frac{60 \times 30}{60 + 30} = \qty{20}{\ohm}$.
\end{solution}

\begin{solution}{exo:g11:circuits-and-power:9}
$\mathcal{E} - 0.50\,r = 5.5$ y $\mathcal{E} - 2.0\,r = 4.0$; restando,
$1.5\,r = 1.5$, luego $r = \qty{1.0}{\ohm}$ y
$\mathcal{E} = \qty{6.0}{V}$. Cortocircuito:
$\mathcal{E}/r = \qty{6.0}{A}$.
\end{solution}

\begin{solution}{exo:g11:circuits-and-power:10}
Hervidor $2200/230 = \qty{9.6}{A}$, tostadora
$1000/230 = \qty{4.3}{A}$: en total \qty{13.9}{A} $< \qty{16}{A}$, el
\omterm{def:g11:circuits-and-power:joule}{fusible} aguanta. Estufa $1500/230 = \qty{6.5}{A}$: en total
\qty{20.4}{A} $> \qty{16}{A}$, se funde (o, lo que es lo mismo,
$\qty{4.7}{kW} > 230 \times 16 \approx \qty{3.7}{kW}$).
\end{solution}

\begin{solution}{exo:g11:circuits-and-power:11}
A \qty{200}{V}: $I = 10^4/200 = \qty{50}{A}$, pérdida
$RI^2 = 2.0 \times 50^2 = \qty{5.0}{kW}$ --- la mitad de lo entregado. A
\qty{4.0}{kV}: $I = \qty{2.5}{A}$, pérdida
$2.0 \times 2.5^2 = \qty{12.5}{W}$. Cociente $400 = 20^2$: veinte veces
la \omterm{def:g11:circuits-and-power:voltage}{tensión} significa veinte veces menos intensidad, y la pérdida va con
$I^2$.
\end{solution}

\begin{solution}{exo:g11:circuits-and-power:12}
Cuatro montajes: las tres en serie, $3 \times 60 = \qty{180}{\ohm}$; las
tres en paralelo, $60/3 = \qty{20}{\ohm}$; una en serie con dos en
paralelo, $60 + 30 = \qty{90}{\ohm}$; una en paralelo con dos en serie,
$\frac{60 \times 120}{180} = \qty{40}{\ohm}$.
\end{solution}

\begin{solution}{exo:g11:circuits-and-power:13}
\emph{1.} $I = \mathcal{E}/(R + r) = 4.5/7.5 = \qty{0.60}{A}$;
$U = RI = \qty{3.6}{V}$; \omterm{def:g10:energy-conservation:power}{potencia} útil $P = UI = \qty{2.2}{W}$; perdida
$rI^2 = 1.5 \times 0.60^2 = \qty{0.54}{W}$;
$\eta = U/\mathcal{E} = 3.6/4.5 = 0.80$.

\emph{2.} Por la pista,
$P = \mathcal{E}^2 \big/ \big[(R - r)^2/R + 4r\big]$; el corchete es
mínimo ($= 4r$) exactamente cuando $R = r$, luego
$P_{\max} = \mathcal{E}^2/(4r) = 4.5^2/6.0 = \qty{3.4}{W}$.
\end{solution}

\begin{solution}{exo:g11:circuits-and-power:14}
Calor: $Q = 1.0 \times 4180 \times (100 - 20) = \qty{3.3e5}{J}$.
Eléctrica: $E = Q/0.90 = \qty{3.7e5}{J}$. Tiempo:
$t = E/P = 3.7\times10^{5}/2200 \approx \qty{170}{s}$, menos de tres
minutos. Coste:
$3.7\times10^{5}/3.6\times10^{6} \approx \qty{0.10}{kW.h}$, unos $2.6$
céntimos.
\end{solution}

\begin{solution}{exo:g11:circuits-and-power:15}
\emph{1.} $q = 60 \times 3600 = \qty{2.2e5}{C}$;
$E \approx \mathcal{E}q = 12.7 \times 2.16\times10^{5} = \qty{2.7e6}{J}
\approx \qty{0.76}{kW.h}$.

\emph{2.} $U = 12.7 - 0.020 \times 150 = \qty{9.7}{V}$; entregada
$UI = \qty{1.5}{kW}$; perdida
$rI^2 = 0.020 \times 150^2 = \qty{450}{W}$.

\emph{3.} Los faros están conectados a los bornes: mientras arranca, la
\omterm{def:g11:circuits-and-power:voltage}{tensión} en bornes cae de $12.7$ a $\qty{9.7}{V}$, así que reciben menos
\omterm{def:g10:energy-conservation:power}{potencia} y se atenúan.
\end{solution}

\begin{solution}{pb:g11:circuits-and-power:1}
\textbf{1.} Cada segundo,
$E_p = mgh = 5.0\times10^{4} \times 9.81 \times 200 = \qty{9.8e7}{J}$:
\omterm{def:g10:energy-conservation:power}{potencia} disponible \qty{98}{MW}.

\textbf{2.} $0.90 \times 98 = \qty{88}{MW}$.

\textbf{3.} $I = P/U = 8.8\times10^{7}/2.0\times10^{4} = \qty{4.4e3}{A}$.

\textbf{4.} $R_\ell I^2 = 32 \times (4.4\times10^{3})^2 \approx
\qty{6.2e8}{W} = \qty{620}{MW}$ --- siete veces lo que produce la
central: transportar a \qty{20}{kV} es imposible.

\textbf{5.} Con $P = UI$ fijo, lo único que encoge $I$ es subir $U$: un
transformador.

\textbf{6.} $I = 8.8\times10^{7}/4.0\times10^{5} = \qty{220}{A}$.

\textbf{7.} $R_\ell I^2 = 32 \times 220^2 \approx \qty{1.6}{MW}$.

\textbf{8.} $1.6/88 \approx 0.018$: un $1.8\,\%$ aproximadamente.

\textbf{9.} Fracción $= R_\ell I^2/P = R_\ell (P/U)^2/P = R_\ell P/U^2
\propto 1/U^2$: de \qty{20}{kV} a \qty{400}{kV} cae en un factor
$20^2 = 400$, y en efecto
$\qty{620}{MW}/400 \approx \qty{1.6}{MW}$.

\textbf{10.} $R_\ell I = 32 \times 220 \approx \qty{7.1}{kV}$: un
$1.8\,\%$ de \qty{400}{kV}, la misma fracción de la pregunta 8, ya que
$R_\ell I/U = R_\ell I^2/UI$.

\textbf{11.} $\qty{1}{kW.h} = \qty{3.6e6}{J}$: el contador multiplica la
\omterm{def:g10:energy-conservation:power}{potencia} consumida por el tiempo que se consume, y va sumando.

\textbf{12.} $I = 900/230 = \qty{3.9}{A}$;
$R = U^2/P = 230^2/900 \approx \qty{59}{\ohm}$.

\textbf{13.} $E = 900 \times 180 = \qty{1.6e5}{J} = \qty{0.045}{kW.h}$:
alrededor de $1.1$ céntimos.

\textbf{14.} La física de $RI^2$ es idéntica; lo único que cambia es el
propósito --- en la línea el calor sobra, y en la tostadora el calor
\emph{es} el producto.

\textbf{15.} $3.9 + 9.6 = \qty{13.5}{A} < \qty{16}{A}$: aguanta. Estufa
$2000/230 = \qty{8.7}{A}$: en total $\qty{22.2}{A} > \qty{16}{A}$, el
\omterm{def:g11:circuits-and-power:joule}{fusible} se funde.

\textbf{16.} $0.05 \times 13.5^2 = \qty{9.1}{W}$:
$9.1/3100 \approx 0.3\,\%$ de la \omterm{def:g10:energy-conservation:power}{potencia} entregada --- despreciable en
cables cortos.

\textbf{17.} $\eta = 0.90 \times 0.99 \times 0.982 \times 0.99 \times
0.997 \approx 0.86$.

\textbf{18.} Un \omterm{def:g10:energy-conservation:kwh}{kilovatio-hora} tostado necesita
$3.6\times10^{6}/0.86 \approx \qty{4.2e6}{J}$ aguas arriba:
$m = E/(gh) = 4.2\times10^{6}/(9.81 \times 200) \approx \qty{2.1e3}{kg}$,
unos \qty{2.1}{m^3} de agua.

\textbf{19.} $R_\ell P/U^2 = 32 \times 8.8\times10^{7}/(2.0\times
10^{4})^2 \approx 7$. Una ``fracción'' mayor que $1$ significa que la
línea disiparía más de lo que transporta: a esa \omterm{def:g11:circuits-and-power:voltage}{tensión} el suministro,
sencillamente, no puede darse.

\textbf{20.} Un \omterm{def:g10:energy-conservation:kwh}{kilovatio-hora} tostado le cuesta a la montaña unas dos
toneladas --- \qty{2.1}{m^3} --- de agua cayendo \qty{200}{m}, y cerca
del $14\,\%$ de la \omterm{def:g10:energy-conservation:energy}{energía} muere por el camino, sobre todo en las turbinas
y los transformadores (menos del $2\,\%$ en \qty{1000}{km} de línea). El
único truco que hace posible el viaje es el transformador: elevar a
\qty{400}{kV} divide la intensidad entre $20$ y la pérdida por \omterm{def:g11:circuits-and-power:joule}{efecto
Joule} entre $400$.
\end{solution}
